\section{R7: Component retrieval at library scale}\label{sec:r7}

\question{R7.Q1}{What abstraction is useful for dependent, universe-polymorphic signatures?}
\status{Use a refinable, constraint-carrying hypergraph; a graph of unrelated argument heads loses essential conjunctions and correlations.}

Begin with a component descriptor containing its telescope, result head,
rigid type-parameter sharing, selected index shapes, universe constraints, and
class prerequisites. The first abstraction may erase most term indices while
retaining facts such as a vector's element type, whether an index is zero or
a successor, and whether two occurrences share the same type parameter. A
later abstraction can retain linear index expressions or a small predicate
set. The purpose is to reject impossible compositions cheaply, not to solve
full dependent unification outside Lean.

A component with two arguments is a hyperedge requiring both arguments, not
two ordinary edges either of which enables its result. If values are reusable,
model that reuse explicitly: a strict token-consuming Petri net without copying
can exclude $x\mapsto(x,x)$. If the language permits unused inputs, do not
silently impose the relevant-typing restriction that every input must be used.
Retrieval and final ranking may prefer relevance without changing inhabitation.

\begin{proposition}[The requirement on a retrieval abstraction]\label{prop:retrieval}
If every well-typed composition in a candidate grammar maps to an abstract
reachability derivation, then abstract unreachability excludes all such
compositions. If a resource-limited abstract search merely fails to find a
path, or if some constructors and recursion templates have no abstract
counterpart, that conclusion does not follow.
\end{proposition}
\begin{proof}
A concrete composition would supply an abstract derivation by hypothesis,
contradicting abstract unreachability. A timeout is not unreachability; a
missing abstract counterpart invalidates the hypothesis.
\end{proof}

TYGAR's abstract typing and refinement framework is the relevant precedent
\cite{hoogleplus}. In Lean, a failed concrete elaboration is not automatically
a proof of untypeability: higher-order unification can be incomplete, and
instance resolution or normalization can time out. Refine from a definite
rigid constructor clash, a checked incompatible index constraint, or another
justified obstruction. Otherwise deprioritize the candidate and retain the
possibility of a different instantiation.

Use universe constraints explicitly or conservatively forget them. Rejecting
an abstraction because a guessed universe assignment failed is unsound when
another assignment could work. Likewise, a path-length cutoff of three is a
bounded retrieval heuristic, not evidence of general uninhabitability. Keep a
fair fallback provider stream if completeness relative to a larger grammar
is claimed. Store indexes by environment fingerprint and declaration identity,
not only a namespace or package version string.

\question{R7.Q2}{Has type-guided abstraction refinement been combined with type classes?}
\status{Yes. Hoogle+ uses dictionary passing directly; Lean adds contextual and dependent-instance complications.}

The Hoogle+ paper's Section~2.5 translates a class requirement such as
\code{Eq a} into an ordinary dictionary argument. A transition requiring
comparison is enabled only if the dictionary type is inhabited. Instance
constructors can themselves provide dictionaries. The answer to the literal
question is therefore already positive, not an unexplored extension of
TYGAR \cite{hoogleplus}.

For Lean, retain a class requirement as an explicit prerequisite in the
component descriptor, and resolve it when enough type information is known.
A dictionary supplied as a local argument is an ordinary available value.
An inferred dictionary must be justified by Lean's instance mechanism or an
explicit construction admitted by the query's grammar.

Memoizing only by ``class and type'' is insufficient. The full instance goal
can contain multiple parameters, indices, output parameters, and universes.
Local instances and their priorities can change the answer. An unresolved
metavariable may later be instantiated in a way that enables a previously
failing search. Key a resolved result by the normalized complete goal,
relevant local-instance context, environment, and assignment fingerprint;
revalidate it on replay. Cache timeout as a resource event, not nonexistence.

Most importantly, computational dictionaries are not proof-irrelevant.
Two \code{Add A} instances can produce different programs; two equality
implementations can yield different observations. Their actual identity and
behavior belong in the candidate and its observational cache. Only a
proposition-valued field enjoys the corresponding proof-irrelevance guarantee.
This is why type-class caching and candidate equivalence must share a semantic
identity policy instead of independently erasing all instance arguments.

\question{R7.Q3}{Can existing premise retrieval support data goals and examples, and where should training data come from?}
\status{Data-goal lookup is already possible; example-conditioned program relevance needs a separate adapter and an appropriate benchmark.}

The pinned Lean library-search source includes a data-goal example,
\code{example : Nat := by exact?}. Thus basic retrieval is not intrinsically
restricted to propositions \cite{librarysearch}. However, accepting a goal
metavariable is not the same as learning how candidate functions transform
example inputs. A premise selector trained on theorem proofs may offer useful
symbol overlap while missing the relevant behavior of a polymorphic function.

A Leant2 retrieval adapter should combine a type-shape query with a separate
behavioral sketch: available input types, observed calls, output constructors,
size changes, and simple relations such as preserving length. The existing
selector can supply one ranked source. A deterministic symbolic index supplies
another. A behavior-aware model, when available, supplies a third. Do not claim
that passing extra prose to a theorem retriever gives a validated
example-conditioned ranking interface.

Recent retrieval work, including LeanHammer and LeanSearch v2, improves
premise retrieval for theorem proving, but its reported evaluations are not
measurements of carrier- or example-conditioned program retrieval
\cite{leanhammer,leansearch}. Hoogle+'s published evaluation uses hundreds of
components, so it should not be cited as direct experimental evidence for
Leant2-scale retrieval across every imported declaration \cite{hoogleplus}.

A useful program corpus can be constructed from small safe definitions in
Init, Std, and selected Mathlib modules. Extract types, explicit component
dependencies from elaborated bodies, and generated observations where the
domains are executable. Train on component use, not merely names occurring in
pretty-printed text. Remove the target definition from the provider set and
remove target-specific equation lemmas or wrappers that trivially reveal its
body. Group near-duplicate algorithms and their lemmas into the same split;
respect import and source-order availability.

Maintain two evaluation tracks: API composition, where existing implementations
may legitimately be used, and body synthesis, where the target and equivalent
wrappers are forbidden. Report retrieval recall of a usable component set and
end-to-end certified synthesis separately. Compiled behavioral probes can
improve ranking, but failure of a standalone component on a whole-function
example cannot refute its usefulness inside a composition.
